\clearpage
\section*{Appendix R\quad Computational Receipts}\label{app:receipts}
\addcontentsline{toc}{section}{Appendix R: Computational Receipts}
\small\noindent Each computation marked \rcptmarker\ in the text is verified by a runnable receipt in the \texttt{receipts/} directory that ships with this work; status \textsf{OK} = verified (traced, run, output matches the claim). This appendix is generated from the receipt index, not maintained by hand.\par\medskip
\begingroup\renewcommand{\arraystretch}{1.25}
\begin{description}
\item[\label{rcpt:P17_power_of_a_point}\texttt{P17\_power\_of\_a\_point}]\hfill\textsf{[OK]}\\
\textit{sec:power} \ (power of a point pow(P)=\textbar{}X\textbar{}\textasciicircum{}2-alpha\textasciicircum{}2 (Steiner/Euclid III.36) = invariant: hinge power=3 exact; 500 random sightlines PX.PY=\textbar{}PO\textbar{}\textasciicircum{}2-alpha\textasciicircum{}2 to 4.4e-16 [borrowed from P17]). 
\emph{Computes:} every-sightline invariant test at machine precision
\ \emph{Run:} \texttt{python3 receipts/P17\_geometric\_core\_paper/P17\_power\_of\_a\_point.py}
\item[\label{rcpt:P17_power_is_null}\texttt{P17\_power\_is\_null}]\hfill\textsf{[OK]}\\
\textit{sec:power prop:tangentnull} \ (tangent to throat is NULL: tangent length=\textbar{}X0\textbar{}=sqrt(R\textasciicircum{}2-alpha\textasciicircum{}2) identity (Euclid=hyperboloid RHS); ds\textasciicircum{}2=-dX0\textasciicircum{}2+\textbar{}dX\textbar{}\textasciicircum{}2=0, max 5.68e-14 over 600 heights (=paper's 5.7e-14); power-of-a-point = null condition, same equation [borrowed from P17]). 
\emph{Computes:} tangent-length identity + ds\textasciicircum{}2=0 at 600 random points
\ \emph{Run:} \texttt{python3 receipts/P17\_geometric\_core\_paper/P17\_power\_is\_null.py}
\item[\label{rcpt:P17_qm_S4_vs_S5}\texttt{P17\_qm\_S4\_vs\_S5}]\hfill\textsf{[OK]}\\
\textit{sec:ledger} \ (su(3) requires so(6)/S\textasciicircum{}5 not so(5)/S\textasciicircum{}4: 8 Gell-Mann generators realified to R\textasciicircum{}6 all antisymmetric (max\textbar{}R+R\textasciicircum{}T\textbar{}=0.0 -> so(6)), faithful on all 6 real dims (minimal real rep R\textasciicircum{}6); GH beta=2pi alpha n-independent -> thermal + gauge Euclideans co-locate on one real form (S\textasciicircum{}5 or S\textasciicircum{}4 equator) [borrowed from P17]). 
\emph{Computes:} realified-generator antisymmetry + faithful-dimension + GH n-independence
\ \emph{Run:} \texttt{python3 receipts/P17\_geometric\_core\_paper/P17\_qm\_S4\_vs\_S5.py}
\end{description}\endgroup
